\begin{abstract}
External-memory systems for language agents must eventually choose which records to
retain. A common reduction assigns one importance score to each record and keeps the
largest scores. Yet the downstream value of a subset may depend on which records are
co-retained. We formulate fixed-reader external memory as a cooperative game and ask
a narrow empirical question: does a low-order interaction model learned from past
subset interventions improve fixed-budget selection and predict lesion damage on
sealed future tasks beyond strong item-only methods? We separate exact M\"obius
coefficients, ordinary Shapley values, Shapley interaction indices, and fitted
polynomial coefficients, and propose an immutable past--freeze--future protocol with
bank-level inference and sign-aware matched lesions. Current experiments do
\emph{not} answer the headline question. Exact and synthetic checks validate the
instrumentation, while real-model diagnostics expose a prior blocker: under missing
evidence, an instruction-tuned reader often emits a correct-looking destination
instead of abstaining, and a formatting intervention can worsen this unsupported
assertion. No tested output contract passed frozen readiness criteria, so prospective
retention remains blocked. This evidence-bounded draft reports the negative
measurement result, corrects an invalid pure-chain argument against Shapley ranking,
and preregisters the tests required to establish---or falsify---incremental value
from memory interactions.
\end{abstract}
